\documentclass[12pt,aspectratio=43]{beamer}

% 自訂字體的封包
\usepackage{fontspec}

% 導入圖形與表格的封包
\usepackage{graphicx}
\usepackage{booktabs}

 % 提供進階數學環境與公式排版
\usepackage{mathtools, amsmath, amsfonts, amsthm, latexsym} 

%% 設定中文字體
\usepackage{xeCJK}
\setCJKmainfont[AutoFakeBold=3.5 , AutoFakeSlant=0.2]{標楷體}
\setmainfont{Times New Roman}
\XeTeXlinebreaklocale "zh"
\XeTeXlinebreakskip = 0pt plus 1pt

% 自訂字體顏色的封包
\usepackage{color} 

% 圖表自動編號的封包
\usepackage{caption} 

% 序列標號
\usepackage{enumerate} 

%% 設定自動編號
\setbeamertemplate{caption}[numbered]

%% 自訂顏色
\definecolor{Myblue}{RGB}{57,108,144}
\definecolor{Default_Blue}{RGB}{52,51,171}

% (動畫)字會若隱若現的顯示
\beamersetuncovermixins{\opaqueness<1->{45}}{\opaqueness<1->{95}} 

% 設定中文的標籤
\renewcommand{\figurename}{圖} 
\renewcommand{\tablename}{表} 

% 排版形式 
\mode<presentation>{
\usetheme{Berlin}
% enumerate 及 item 的形狀
\setbeamertemplate{itemize items}[triangle]
% 調整外框形式的字體大小
\setbeamerfont{headline}{size=\scriptsize}
\setbeamerfont{footline}{size=\scriptsize}
% 取消右下方的跳轉工具列
\setbeamertemplate{navigation symbols}{} 
% 自定義2：清除外框下緣但僅出頁碼
\setbeamertemplate{footline}[page number] 
% 顏色主題
\usecolortheme{dove}
% 設定頁面邊界
\setbeamersize{text margin left=0.6cm, text margin right=0.6cm}
\special{papersize=\the\paperwidth,\the\paperheight}
\providecommand{\tabularnewline}{\\}
}

\title{Meeting 1015}
\author{黃丞暘}
\institute{國立東華大學應數統計所}
\date{\today}

\begin{document}

% 封面
\begin{frame}
  \titlepage
\end{frame}

% --- 目錄 ---
\begin{frame}{\textbf{目錄}}
  \tableofcontents
\end{frame}

%	--- 文獻回顧I ---
\section{文獻回顧I}
\linespread{1}

\begin{frame}[<+>]{\textbf{文獻回顧I}}
\linespread{1.5}
	\begin{itemize}
		\item Lai ( 2002 )
		\begin{enumerate}
			\item 方法
                \begin{enumerate}
                    \item Data preprocessing
                    \item KNN-imputation
                    \item Random forest features selection
                    \item imbalance - SMOTE to 1:1
                      \only<7>{%
                        \begin{center}
                          \includegraphics[width=0.65\textwidth]{C:/Users/USER/Documents/GitHub/psychic-spoon/0 Meeting/20251015_Meeting/SMOTE_AFTER.png}
                          \captionof{figure}{SMOTE after}
                        \end{center}
                      }
                    \item<8-> 0.8 training, 0.2 testing by decision tree
                \end{enumerate}
    \item<9-> 結果
                \begin{enumerate}
                    \item<10-> Best: ACC 86.94\%, FAR 9.22\%
                    \item<11-> 反查 Fail 節點
                \end{enumerate}
    \end{enumerate}
  \end{itemize} 
\end{frame}

\section{文獻回顧I}

\begin{frame}[<+>]{\textbf{文獻回顧I}}
\linespread{1.5}

\begin{itemize}
  \item Lai (2002)
  \begin{enumerate}
    \item 方法
      \begin{enumerate}
        \item Data preprocessing
        \item KNN-imputation
        \item Random forest features selection
        \item imbalance - SMOTE to 1:1

        % ======= 第二階段：只顯示圖片 =======
        \only<6>{%
          \begin{center}
            \includegraphics[width=0.65\textwidth]{C:/Users/USER/Documents/GitHub/psychic-spoon/0 Meeting/20251015_Meeting/SMOTE_AFTER.png}
            \captionof{figure}{SMOTE after}
          \end{center}
        }

        % ======= 第三階段：回到原內容 =======
        \item<7-> 0.8 training, 0.2 testing by decision tree
      \end{enumerate}

    \item<8-> 結果
      \begin{enumerate}
        \item<9-> Best: ACC 86.94\%, FAR 9.22\%
        \item<10-> 反查 Fail 節點
      \end{enumerate}
  \end{enumerate}
\end{itemize}
\end{frame}

% --- Frame 1: 圖片出現前的內容 ---
\begin{frame}[<+>]{\textbf{文獻回顧I}}
\linespread{1.5}
  \begin{itemize}
    \item Lai (2002)
    \begin{enumerate}
      \item 方法
      \begin{enumerate}
        \item Data preprocessing
        \item KNN-imputation
        \item Random forest features selection
        \item imbalance - SMOTE to 1:1
      \end{enumerate}
    \end{enumerate}
  \end{itemize}
\end{frame}

% --- Frame 2: 獨立顯示圖片 ---
\begin{frame}{\textbf{文獻回顧I} (SMOTE 示意圖)}
  \begin{figure}
    \centering
    \includegraphics[width=0.8\textwidth]{C:/Users/USER/Documents/GitHub/psychic-spoon/0 Meeting/20251015_Meeting/SMOTE_AFTER.png}
    \caption{SMOTE after}
  \end{figure}
\end{frame}

% --- Frame 3: 接續圖片之後的內容 ---
\begin{frame}[<+>]{\textbf{文獻回顧I}}
\linespread{1.5}
  \begin{itemize}
    \item Lai (2002)
    \begin{enumerate}
      \item 方法
      \begin{enumerate}
        % --- 以下項目沒有 < > 標籤 ---
        % 這代表它們在這一頁出現時，會立即顯示 (通常是淡化狀態)
        % 作為動畫的背景/上下文
        \item Data preprocessing
        \item KNN-imputation
        \item Random forest features selection
        \item imbalance - SMOTE to 1:1

        % --- 動畫從這裡開始 ---
        % 第一個 <+-> 標籤，代表按下「下一步」後，才會出現的第一個新項目
        \item<+-> 0.8 training, 0.2 testing by decision tree
      \end{enumerate}
      % 後續的動畫會依序播放
      \item<+-> 結果
      \begin{enumerate}
        \item<+-> Best: ACC 86.94\%, FAR 9.22\%
        \item<+-> 反查 Fail 節點
      \end{enumerate}
    \end{enumerate}
  \end{itemize}
\end{frame}

\end{document}
